%% ============================================================
%% Section 03d — Parse-Time Enforcement (Parser/Planner Gate)
%% Author:     Barbara — Lead / Architect
%% Date:       2026-06-05T00:14:04-04:00
%% Stream:     F — Parser/Planner Gate Specification (Phase 1)
%% Status:     Normative draft — Phase 2 runtime implementors are
%%             bound by this section.
%% References:
%%   design/gert/grammar/gxl.ebnf   (GXL error codes, §6)
%%   design/gert/grammar/gis.ebnf   (GIS error codes, §5)
%%   design/gert/grammar/gcp.ebnf   (GCP error codes, §7)
%%   design/gert/expression-language-proposal.md §4
%%   design/web-platform/c-sharp-governance-parity.md §4
%%   .squad/decisions.md (OI-GIS-01 ratification, OI-GCP-02 ratification)
%% ============================================================

\section{Parse-Time Enforcement}
\label{sec:parse-gate}

\subsection{Governance Rationale}
\label{sec:parse-gate:rationale}

The words \textbf{MUST}, \textbf{MUST NOT}, \textbf{REQUIRED}, \textbf{SHALL},
\textbf{SHALL NOT}, \textbf{SHOULD}, \textbf{SHOULD NOT}, \textbf{RECOMMENDED},
\textbf{MAY}, and \textbf{OPTIONAL} in this section are to be interpreted as
described in RFC~2119.

GERT's value proposition rests on three pillars: governance by construction,
audit trail integrity, and cross-runtime parity.  None of these properties
survives a runtime that begins executing steps before every expression in the
runbook has been inspected against the canonical grammars.

\begin{description}
  \item[Cross-runtime parity.]
    A runbook that runs on the Go reference implementation
    \emph{must} produce identical observable behaviour on the C\# runtime and
    on any future runtime.  That guarantee is only achievable if every runtime
    parses the same canonical grammars---\texttt{gxl.ebnf},
    \texttt{gis.ebnf}, \texttt{gcp.ebnf}---before committing to an execution
    plan.  Lazy or runtime-only parsing creates a window in which one
    implementation accepts an expression another rejects, silently diverging
    audit trails.

  \item[Audit trail integrity.]
    The JSONL trace is the authoritative, tamper-evident record of every
    runbook execution.  Trace integrity requires that the set of expressions
    executed is \emph{exactly} the set that was validated.  An expression
    discovered to be invalid mid-execution cannot be cleanly rolled back; the
    run is in an undefined state and the trace is untrustworthy.

  \item[Replay determinism.]
    GERT supports evidence replay.  A trace recorded in environment~A must
    produce the same step sequence when replayed in environment~B.  Replay
    determinism requires that expression syntax is fixed at plan time, not
    discovered dynamically.  If an expression is invalid, the replay must fail
    at the plan gate, not partway through re-execution.
\end{description}

For these reasons, parse-time enforcement is \textbf{non-negotiable} and is the
\emph{single} governance checkpoint for expression validity.  There are no
alternative paths, no fast paths, and no deferred-validation modes.

% ─────────────────────────────────────────────────────────────────────────────
\subsection{The Parse-Time Contract}
\label{sec:parse-gate:contract}
% ─────────────────────────────────────────────────────────────────────────────

\paragraph{Definition: Validated Runbook.}
A runbook is \emph{validated} if and only if:

\begin{enumerate}
  \item Every field that carries a GXL boolean expression---\texttt{when},
    \texttt{condition}, \texttt{until}, \texttt{iterate.until},
    \texttt{branches[\,].condition}, \texttt{include.when},
    \texttt{collector.fields[\,].when}---has been successfully parsed against
    the GXL grammar (\texttt{gxl.ebnf}, Version~1.0.0-draft or later) with no
    parse errors.

  \item Every field that carries a GIS template string---\texttt{title},
    \texttt{args}, \texttt{stdin}, \texttt{tool.argv[\,]},
    \texttt{display.content}, \texttt{assert[\,].subject},
    \texttt{assert[\,].expected}, artifact paths, \texttt{include.with}
    values---has been successfully parsed against the GIS grammar
    (\texttt{gis.ebnf}, Version~1.0.0-draft or later) with no parse
    \emph{errors} (GIS-PARSE-002 warnings do not block validation; all other
    GIS error codes do).

  \item Every capture path in every \texttt{capture:} map has been
    successfully parsed against the GCP grammar (\texttt{gcp.ebnf},
    Version~1.0.0-draft or later) with no parse errors.

  \item Every semantic precondition checked by the Planner
    (Section~\ref{sec:parse-gate:validated-plan}) has passed.
\end{enumerate}

\paragraph{MUST: No step execution from an unvalidated plan.}
\texttt{RunHandle.Next()} (Go) and \texttt{IRunHandle.NextAsync()} (C\#)---and
their equivalent in any future runtime---\textbf{MUST NOT} execute any step
from a plan that has not passed the full validation pipeline described in this
section.  A runtime that receives anything other than a
\texttt{ValidatedPlan}~(\S\ref{sec:parse-gate:validated-plan}) as input to its
run-start function \textbf{MUST} refuse with a hard error before any step
dispatches.

\paragraph{MUST: Validation is performed once, at plan time.}
Validation \textbf{MUST} be performed in full before the first step dispatches.
Lazy parsing---deferring expression parsing until the step that contains the
expression is about to execute---is \textbf{MUST NOT}.  The rationale is
audit trail integrity: a mid-execution parse failure leaves the run in an
undefined state whose trace cannot be replayed.

The validation pipeline sequence is:

\begin{enumerate}
  \item \textbf{YAML Load:} structural parsing of the runbook YAML; JSON Schema
    / structural validation.
  \item \textbf{GXL Parse:} all boolean-expression fields parsed by the GXL
    parser.
  \item \textbf{GIS Parse:} all string-template fields parsed by the GIS
    parser.
  \item \textbf{GCP Parse:} all capture paths parsed by the GCP parser.
  \item \textbf{Semantic validation (Planner):} variable scope analysis,
    undefined-capture detection, type inference, subtree-default prohibition
    (\texttt{GCP-DEFAULT-SUBTREE}), grammar-version recording.
\end{enumerate}

If \emph{any} step fails, the pipeline \textbf{MUST} stop, accumulate all
errors from the failed stage, and return them to the caller.  No
\texttt{ValidatedPlan} is produced.  No \texttt{RunHandle} is created.

\paragraph{MUST: Validation result recorded in trace.}
When a runbook is successfully validated, the runtime \textbf{MUST} emit the
event \texttt{plan.validated} as the \emph{first} event after plan load and
before any step event.  The event payload \textbf{MUST} include at minimum:

\begin{itemize}
  \item \texttt{runbook\_id} and \texttt{runbook\_hash} (content hash of the
    source file).
  \item \texttt{grammar\_versions} as recorded in the
    \texttt{ValidatedPlan}~(\S\ref{sec:parse-gate:validated-plan}).
  \item \texttt{validated\_at} timestamp.
  \item \texttt{expression\_count} summary (GXL, GIS, GCP expression counts).
\end{itemize}

The trace event schema for \texttt{plan.validated} is defined in
\texttt{design/gert/sections/07-runtime-events.tex}.  Implementations \textbf{MUST}
reference that schema for the full event envelope.

% ─────────────────────────────────────────────────────────────────────────────
\subsection{The \texttt{ValidatedPlan} Type Contract}
\label{sec:parse-gate:validated-plan}
% ─────────────────────────────────────────────────────────────────────────────

A \texttt{ValidatedPlan} is the \emph{only} value accepted by the run-start
function of any GERT runtime.  It is a closed, opaque product of the validation
pipeline; it cannot be constructed by caller code directly.

\paragraph{Normative field specification.}
Every \texttt{ValidatedPlan} value \textbf{MUST} expose the following fields
(field names \textbf{MAY} follow language naming conventions---\texttt{camelCase}
for C\#, \texttt{snake\_case} for Go---but the semantics and types \textbf{MUST}
be equivalent):

\begin{center}
\begin{tabular}{lll}
\hline
\textbf{Field} & \textbf{Type} & \textbf{Description} \\
\hline
\texttt{runbook\_id}      & string    & Unique identifier of the runbook \\
\texttt{runbook\_hash}    & string    & SHA-256 hex digest of the source file \\
\texttt{grammar\_versions}& struct    & Grammar versions used during validation \\
\quad\texttt{.gxl}        & string    & e.g., \texttt{"1.0.0-draft"} \\
\quad\texttt{.gis}        & string    & e.g., \texttt{"1.0.0-draft"} \\
\quad\texttt{.gcp}        & string    & e.g., \texttt{"1.0.0-draft"} \\
\texttt{validated\_at}    & timestamp & RFC~3339 UTC timestamp of validation \\
\texttt{expression\_count}& struct    & Count of successfully parsed expressions \\
\quad\texttt{.gxl}        & int       & Number of GXL expression fields \\
\quad\texttt{.gis}        & int       & Number of GIS template fields \\
\quad\texttt{.gcp}        & int       & Number of GCP capture paths \\
\texttt{errors}           & list      & \textbf{MUST} be empty (see below) \\
\hline
\end{tabular}
\end{center}

The \texttt{errors} field \textbf{MUST} be empty for any \texttt{ValidatedPlan}
that is accepted by the runtime.  A plan with a non-empty \texttt{errors} list
is an invalid plan; it \textbf{MUST} be rejected and \textbf{MUST NOT} be
stored in any persistent plan store.  The type invariant is: if
\texttt{errors} is non-empty, the value is not a \texttt{ValidatedPlan}---it
is an error report.

\paragraph{Language-specific implementation notes.}
Go implementations \textbf{MUST} represent \texttt{ValidatedPlan} as an
unexported struct with no exported constructor other than the planner's
\texttt{Plan()} function.  C\# implementations \textbf{MUST} represent
\texttt{ValidatedPlan} as an \texttt{internal sealed record} accessible only
via \texttt{IRunbookPlanner.PlanAsync()}.  No adapter, controller, or external
caller \textbf{MAY} construct a \texttt{ValidatedPlan} by any means other than
the canonical planner.  See
\texttt{design/web-platform/c-sharp-governance-parity.md~\S1.2} for the
normative C\# pipeline class mapping.

% ─────────────────────────────────────────────────────────────────────────────
\subsection{Error Code Catalog}
\label{sec:parse-gate:error-codes}
% ─────────────────────────────────────────────────────────────────────────────

This subsection is the \textbf{normative master error-code list} for all three
grammar parsers and the parse gate.  Phase~2 implementors \textbf{MUST} emit
exactly these codes.  Error codes from individual grammar files are reproduced
here for completeness; the grammar files remain authoritative for parse-stage
codes.

\subsubsection{GXL Parser Error Codes}
\label{sec:parse-gate:error-codes:gxl}

Source: \texttt{design/gert/grammar/gxl.ebnf~\S6}.
Raised by: GXL Parser (plan time unless noted).

\begin{center}\small
\begin{tabular}{p{3.5cm}p{1.8cm}p{5.5cm}p{3.5cm}}
\hline
\textbf{Code} & \textbf{Raised By} & \textbf{Condition} & \textbf{Remediation} \\
\hline
\texttt{GXL-PARSE-001} & Parser & Unexpected token & Check expression syntax \\
\texttt{GXL-PARSE-002} & Parser & Unterminated string literal & Close the quote \\
\texttt{GXL-PARSE-003} & Parser & Invalid number literal (leading zeros, NaN, Infinity) & Use finite decimal form \\
\texttt{GXL-PARSE-004} & Parser & Invalid escape sequence & Use supported escapes (\textbackslash\textbackslash, \textbackslash n, \textbackslash t, \textbackslash r, \textbackslash u, \textbackslash U) \\
\texttt{GXL-PARSE-005} & Parser & Unknown namespace (\texttt{math.*} in v1) & Remove \texttt{math.*} calls \\
\texttt{GXL-PARSE-006} & Parser & Unknown namespace method & Check stdlib reference \\
\texttt{GXL-PARSE-007} & Parser & Forbidden syntax: \texttt{\&\&}, \texttt{||}, \texttt{!}, \texttt{|}, ternary, infix \texttt{contains}/\texttt{startsWith}/\texttt{matches}, assignment, string concat, array/object construction, negative index & Use GXL equivalents \\
\texttt{GXL-PARSE-008} & Parser & Unmatched parenthesis & Balance parentheses \\
\texttt{GXL-PARSE-009} & Parser & Chained comparison (\texttt{a < b < c}) & Use \texttt{a < b and b < c} \\
\texttt{GXL-PARSE-010} & Parser & Keyword used as identifier & Rename the variable \\
\texttt{GXL-TYPE-001}  & Parser/Planner & Type mismatch in comparison & Use same-type operands \\
\texttt{GXL-TYPE-002}  & Parser/Planner & Non-bool in boolean position & Ensure expression resolves to bool \\
\texttt{GXL-TYPE-003}  & Evaluator & Argument type mismatch / arithmetic on non-number & Check argument types \\
\texttt{GXL-TYPE-004}  & Parser & Wrong argument count & Check function signature \\
\texttt{GXL-PATH-001}  & Evaluator & Variable or path segment not found & Declare variable or use \texttt{when:} guard \\
\texttt{GXL-PATH-002}  & Evaluator & Array index out of bounds & Guard with \texttt{len()} check \\
\texttt{GXL-PATH-003}  & Evaluator & Index on non-array & Check capture type \\
\texttt{GXL-EVAL-001}  & Evaluator & Variable resolution failure (catch-all) & Inspect context map \\
\texttt{GXL-EVAL-002}  & Evaluator & Division/modulo by zero & Guard divisor \\
\texttt{GXL-EVAL-003}  & Evaluator & Invalid RE2 regex pattern & Use RE2 syntax \\
\texttt{GXL-EVAL-004}  & Evaluator & Null in ordered comparison & Use \texttt{!= null} guard \\
\hline
\end{tabular}
\end{center}

\subsubsection{GIS Parser Error Codes}
\label{sec:parse-gate:error-codes:gis}

Source: \texttt{design/gert/grammar/gis.ebnf~\S5}.
Raised by: GIS Parser (plan time) or GIS Evaluator (runtime).

\begin{center}\small
\begin{tabular}{p{3.5cm}p{1.8cm}p{5.5cm}p{3.5cm}}
\hline
\textbf{Code} & \textbf{Raised By} & \textbf{Condition} & \textbf{Remediation} \\
\hline
\texttt{GIS-PARSE-001}    & Parser & Unclosed interpolation block (\texttt{\$\{} without \texttt{\}}) & Add closing \texttt{\}} \\
\texttt{GIS-PARSE-002}    & Parser & Unknown escape sequence (WARNING only, not hard error) & Use supported escapes or literal \textbackslash\$\{ \\
\texttt{GIS-PARSE-003}    & Parser & Invalid GXL expression inside \texttt{\$\{...\}} & Fix embedded GXL; sub-error code included \\
\texttt{GIS-PARSE-004}    & Parser & Empty interpolation block (\texttt{\$\{\}}) & Provide variable name \\
\texttt{GIS-PATH-MISSING} & Evaluator & GXL expression raised \texttt{GXL-PATH-001} & Declare variable or use \texttt{when:} guard \\
\texttt{GIS-TYPE-001}     & Evaluator & NaN or Infinity in string coercion & Ensure capture yields finite number \\
\texttt{GIS-TYPE-002}     & Evaluator & GXL type error propagated from embedded expression & Fix GXL expression type \\
\hline
\end{tabular}
\end{center}

\subsubsection{GCP Parser Error Codes}
\label{sec:parse-gate:error-codes:gcp}

Source: \texttt{design/gert/grammar/gcp.ebnf~\S7}.
Raised by: GCP Parser (plan time) or GCP Evaluator (runtime).

\begin{center}\small
\begin{tabular}{p{3.5cm}p{1.8cm}p{5.5cm}p{3.5cm}}
\hline
\textbf{Code} & \textbf{Raised By} & \textbf{Condition} & \textbf{Remediation} \\
\hline
\texttt{GCP-PARSE-001}      & Parser   & Unknown source prefix (e.g., \texttt{xml.*}) & Use a supported prefix \\
\texttt{GCP-PARSE-002}      & Parser   & Invalid path suffix (e.g., \texttt{exit\_code.extra}) & Remove suffix \\
\texttt{GCP-PARSE-003}      & Parser   & Negative array index & Use non-negative index \\
\texttt{GCP-PARSE-004}      & Parser   & Empty path segment (double dot or trailing dot) & Remove extra dot \\
\texttt{GCP-PARSE-005}      & Parser   & Invalid header name & Use lowercase alphanumeric with hyphens \\
\texttt{GCP-PARSE-006}      & Planner  & \texttt{step.\{id\}.*} references undefined step ID & Correct step ID \\
\texttt{GCP-RESOLVE-001}    & Evaluator & Referenced step not yet executed or was skipped & Re-order steps or guard with \texttt{when:} \\
\texttt{GCP-RESOLVE-002}    & Evaluator & Path segment not found in object & Check step output structure \\
\texttt{GCP-RESOLVE-003}    & Evaluator & Array index out of bounds / index on non-array & Guard with \texttt{len()} \\
\texttt{GCP-RESOLVE-004}    & Evaluator & Output not parseable as declared format (json/yaml) & Check step output format \\
\texttt{GCP-RESOLVE-005}    & Evaluator & Header absent (returns null, not error, by design) & Use \texttt{capture.default:} if required \\
\texttt{GCP-DEFAULT-SUBTREE}& Planner  & \texttt{capture.default:} declared for subtree capture & Use \texttt{when: var != null} guard \\
\texttt{GCP-TYPE-001}       & Planner  & Default value type incompatible with capture path & Match default type to capture type \\
\hline
\end{tabular}
\end{center}

\subsubsection{Parse-Gate Error Codes (Plan Level)}
\label{sec:parse-gate:error-codes:plan}

The following codes are raised by the \textbf{Planner} (semantic validation
layer) after all three parsers have succeeded, or by the gate layer itself.
They are \emph{new in Stream~F} and do not originate from any grammar file.
They are proposed in
\texttt{.squad/decisions/inbox/barbara-stream-f-error-codes.md} and included
here as normative.

\begin{center}\small
\begin{tabular}{p{2.8cm}p{1.8cm}p{6.0cm}p{3.7cm}}
\hline
\textbf{Code} & \textbf{Raised By} & \textbf{Condition} & \textbf{Remediation} \\
\hline
\texttt{PLAN-001} & Parser / Planner & Runbook contains a field with an unparseable expression (wrapper; includes sub-error from grammar) & Fix the flagged expression field \\
\texttt{PLAN-002} & Planner & Runbook references an undefined capture variable in an expression (variable not produced by any reachable predecessor step) & Ensure the variable is captured before use \\
\texttt{PLAN-003} & Planner & \texttt{capture.default:} declared on a subtree capture (alias for \texttt{GCP-DEFAULT-SUBTREE}; retained for plan-level reporting) & Use \texttt{when: var != null} guard \\
\texttt{PLAN-004} & Parser  & Forbidden Go-template syntax \texttt{\{\{\ \}\}} detected in a GIS field & Replace with \texttt{\$\{...\}} interpolation \\
\texttt{PLAN-005} & Parser  & Forbidden infix operator (\texttt{\&\&}, \texttt{||}, \texttt{!}, \texttt{contains} used as infix) detected in a GXL field & Replace with \texttt{and}, \texttt{or}, \texttt{not}, \texttt{str.contains()} \\
\texttt{PLAN-006} & Parser  & Forbidden pipe expression in a GCP capture path (\texttt{|\ length} or similar) & Replace with \texttt{len()} in a GXL expression after scalar capture \\
\texttt{PLAN-007} & Planner & Grammar version mismatch: recorded \texttt{grammar\_versions} in a stored plan are incompatible with the validating runtime's grammar versions & Re-validate the runbook against the current grammar \\
\texttt{PLAN-008} & Planner & Keyword used as a capture variable name (e.g., \texttt{capture: and: ...}) & Rename the capture variable (\texttt{GXL-PARSE-010} for GXL context) \\
\texttt{PLAN-009} & Planner & Runbook references a \texttt{step.X.*} capture path where step~X is not statically reachable from the declaring step (conditional branch exclusion) & Re-order steps or use \texttt{when:} guard \\
\hline
\end{tabular}
\end{center}

\noindent\textit{Note:} \texttt{PLAN-004}, \texttt{PLAN-005}, and
\texttt{PLAN-006} are migration-era codes designed to catch runbooks that have
not yet been processed by the Stream~E migration tooling.  Once all fixtures
are migrated, these codes remain active as defensive guards against accidental
re-introduction of deprecated syntax.

% ─────────────────────────────────────────────────────────────────────────────
\subsection{Error Message Format}
\label{sec:parse-gate:error-format}
% ─────────────────────────────────────────────────────────────────────────────

Every error emitted by the validation pipeline \textbf{MUST} conform to the
following structured format.  Implementations \textbf{MUST NOT} emit
unstructured string messages as the primary error representation.

\begin{center}\small
\begin{tabular}{p{2.8cm}p{2.2cm}p{7.5cm}}
\hline
\textbf{Field} & \textbf{Type} & \textbf{Description} \\
\hline
\texttt{code}     & string  & Error code from the catalog in~\S\ref{sec:parse-gate:error-codes} \\
\texttt{message}  & string  & Human-readable description in English \\
\texttt{location} & struct  & Source location of the offending token or field \\
\quad\texttt{.file}       & string & Runbook or tool file path (relative to runbook root) \\
\quad\texttt{.line}       & int    & 1-based line number \\
\quad\texttt{.column}     & int    & 1-based column (code-point offset, not byte offset) \\
\quad\texttt{.span\_length} & int  & Length in code points of the highlighted span \\
\texttt{snippet}    & string & The offending source slice, extracted verbatim \\
\texttt{suggestion} & string & OPTIONAL machine-generated remediation hint \\
\hline
\end{tabular}
\end{center}

\paragraph{Example (PLAN-004, forbidden Go-template syntax):}

\begin{verbatim}
{
  "code": "PLAN-004",
  "message": "Go-template syntax '{{ }}' is not supported in GIS fields.
               Use '${...}' interpolation instead.",
  "location": { "file": "schema.yaml", "line": 66, "column": 12,
                "span_length": 13 },
  "snippet": "{{ .pod_name }}",
  "suggestion": "Replace '{{ .pod_name }}' with '${pod_name}'"
}
\end{verbatim}

\paragraph{Example (PLAN-005, forbidden infix operator):}

\begin{verbatim}
{
  "code": "PLAN-005",
  "message": "Infix '||' is not a valid GXL operator.",
  "location": { "file": "schema.yaml", "line": 152, "column": 28,
                "span_length": 2 },
  "snippet": "error_rate > 0.05 || p95_latency > 0.5",
  "suggestion": "Replace '||' with 'or'"
}
\end{verbatim}

When errors are recorded in the execution trace (for runs that fail at the
plan gate before any steps execute), they \textbf{MUST} be embedded in a
\texttt{plan.failed} trace event whose schema is defined in
\texttt{design/gert/sections/07-runtime-events.tex}.  Each error in the array
\textbf{MUST} use the structured format above.

% ─────────────────────────────────────────────────────────────────────────────
\subsection{No-Bypass Guarantees}
\label{sec:parse-gate:no-bypass}
% ─────────────────────────────────────────────────────────────────────────────

The parse gate is the \textbf{single} checkpoint for expression validity.
There are no alternative paths, no opt-out flags, no ``fast path'' modes that
skip validation.  This is enforced structurally, not by convention.

\paragraph{Go runtime.}
The Go runtime \textbf{MUST} expose no exported function that constructs a
\texttt{RunHandle} from anything other than a \texttt{*ValidatedPlan} returned
by the canonical planner.  The type system enforces this: \texttt{ValidatedPlan}
is an unexported type producible only by \texttt{Planner.Plan()}.  No package
outside the planner package may forge a \texttt{ValidatedPlan}.

\paragraph{C\# runtime (Go-equivalent enforcement).}
The C\# runtime \textbf{MUST} enforce the same constraint via the
\texttt{internal sealed} pattern documented in
\texttt{design/web-platform/c-sharp-governance-parity.md~\S4.1}:

\begin{itemize}
  \item \texttt{ValidatedPlan} is an \texttt{internal sealed record} in
    assembly \texttt{Gert.Runtime.Core}.
  \item \texttt{IRunHandle} is the only public execution surface.
  \item \texttt{IGertRuntime.StartAsync()} requires a \texttt{ValidatedPlan};
    no overload exists that accepts a raw \texttt{ParsedRunbook} or an
    unvalidated schema.
  \item \texttt{IStepRunner} and all concrete runners remain
    \texttt{internal sealed}---invisible to adapter code.
\end{itemize}

\paragraph{Roslyn analyzers (\texttt{GERT0001} / \texttt{GERT0002}).}
The existing Roslyn analyzers defined in
\texttt{design/web-platform/c-sharp-governance-parity.md} enforce
\texttt{IStepRunner} inaccessibility and RE2 regex compliance.  These are
necessary but not sufficient for parse-gate enforcement.  Every runtime
\textbf{MUST} additionally ensure:

\begin{enumerate}
  \item No adapter assembly has \texttt{InternalsVisibleTo} access to
    \texttt{Gert.Runtime.Core}.
  \item CI build verification tests \textbf{MUST} assert that
    \texttt{IGertRuntime.StartAsync} is the only code path that produces
    a live \texttt{IRunHandle}, and that it is unreachable without a
    \texttt{ValidatedPlan} argument.
  \item Equivalent static analysis \textbf{MUST} exist for the Go runtime:
    no Go build tag or conditional compilation path may bypass the
    \texttt{ValidatedPlan} type constraint.
\end{enumerate}

\noindent
The parse gate is not a policy choice.  It is a structural property of every
conformant GERT runtime.

% ─────────────────────────────────────────────────────────────────────────────
\subsection{Grammar Version Pinning}
\label{sec:parse-gate:version-pinning}
% ─────────────────────────────────────────────────────────────────────────────

Each grammar file declares a \texttt{Version} header in its first comment
block.  At the time of this writing the canonical versions are:

\begin{itemize}
  \item \texttt{gxl.ebnf}: \texttt{1.0.0-draft}
  \item \texttt{gis.ebnf}: \texttt{1.0.0-draft}
  \item \texttt{gcp.ebnf}: \texttt{1.0.0-draft}
\end{itemize}

\paragraph{MUST: Record grammar versions in \texttt{ValidatedPlan}.}
The Planner \textbf{MUST} read the \texttt{Version} header from each grammar
file (or from the compiled parser version constant that corresponds to a
specific grammar file revision) and record the three versions in the
\texttt{grammar\_versions} field of the produced \texttt{ValidatedPlan}.

\paragraph{MUST: Reject on version mismatch.}
When a runtime loads a \texttt{ValidatedPlan} from persistent storage (e.g.,
re-validation or replay), it \textbf{MUST} compare the recorded
\texttt{grammar\_versions} against its own compiled grammar versions.  If any
version is incompatible (different major version, or a minor version with a
documented breaking change), the runtime \textbf{MUST} refuse to execute and
\textbf{MUST} raise \texttt{PLAN-007}.

\paragraph{Compatibility policy.}
The draft-to-release transition from \texttt{1.0.0-draft} to
\texttt{1.0.0} is a breaking change requiring re-validation.  Within a stable
major version, backward-compatible additions (new error codes, new optional
grammar productions) \textbf{SHOULD} increment the minor version; plans
validated against an older minor version \textbf{SHOULD} be re-validated as a
best practice but \textbf{MAY} be accepted if the runtime can confirm no
breaking production changes affect the stored plan.

\paragraph{Open governance question.}
What happens to in-flight executions when grammar versions are upgraded
mid-run?  This scenario---a run that starts on grammar \texttt{1.0.0}
and is mid-checkpoint when the runtime upgrades to \texttt{1.1.0}---is
not resolved in Phase~1.  It is surfaced as a Phase~2 governance decision
(\S\ref{sec:parse-gate:open-issues}, OPQ-GATE-01).

% ─────────────────────────────────────────────────────────────────────────────
\subsection{Open Issues for Phase~2}
\label{sec:parse-gate:open-issues}
% ─────────────────────────────────────────────────────────────────────────────

The following issues were surfaced during Stream~F authoring.  They do not
block Phase~1 deliverables but \textbf{MUST} be resolved before Phase~2
runtime implementation begins.  Each is assigned an identifier for tracking.

\begin{description}
  \item[\texttt{OPQ-GATE-01} — In-flight grammar version upgrade.]
    What is the correct behaviour when the runtime's grammar version changes
    while a run is mid-execution (e.g., mid-approval-gate checkpoint)?  Two
    options:

    \textbf{Option~A (strict):} The run is suspended at the next
    \texttt{RunHandle.Next()} call, \texttt{PLAN-007} is raised, and the run
    is marked \texttt{VersionMismatch}.  The operator must re-validate and
    replay or cancel.

    \textbf{Option~B (sticky):} The run retains the grammar version it was
    validated against for its lifetime; grammar upgrades only affect new runs.
    The runtime maintains a version registry mapping run~IDs to grammar
    versions.

    \textit{Trade-off:} Option~A is safer (eliminates mixed-version execution)
    but operationally disruptive for long-running approval flows.  Option~B is
    operationally friendlier but requires version registry infrastructure and
    complicates the conformance corpus.

    \textbf{Recommendation:} Surface to Germán for decision before Phase~2
    begins.  Memo at \texttt{.squad/decisions/inbox/barbara-stream-f-complete.md}.

  \item[\texttt{OPQ-GATE-02} — Statically reachable step set definition
    (OI-GCP-03).]
    \texttt{GCP-PARSE-006} and \texttt{PLAN-009} require the Planner to know
    which steps are statically reachable.  For runbooks with conditional
    branches (\texttt{branch} steps) or optional includes
    (\texttt{include.when:}), the reachable step set is not a single sequence
    but a DAG.  The Planner \textbf{MUST} define ``statically reachable'' for
    the purposes of cross-step capture validation (\texttt{step.\{id\}.*}).
    Phase~2 must specify the reachability algorithm (conservative over-
    approximation recommended).

  \item[\texttt{OPQ-GATE-03} — GIS-PARSE-004 empty interpolation as
    error vs.\ PLAN-001.]
    \texttt{GIS-PARSE-004} (empty interpolation block) is currently a pure
    GIS parse error.  Should it also produce \texttt{PLAN-001} as a wrapper?
    The error catalog consolidation in~\S\ref{sec:parse-gate:error-codes:plan}
    treats \texttt{PLAN-001} as the top-level wrapper for any parse failure.
    Phase~2 must clarify the aggregation strategy for multi-error plans.

  \item[\texttt{OPQ-GATE-04} — GIS-PARSE-002 warning and
    \texttt{plan.validated} event.]
    \texttt{GIS-PARSE-002} (unknown escape sequence) is a warning, not a hard
    error.  The \texttt{plan.validated} event does not currently carry a
    \texttt{warnings} list.  Phase~2 must decide whether warnings are
    included in the \texttt{plan.validated} event and how they appear in the
    JSONL trace.

  \item[\texttt{OPQ-GATE-05} — Plan storage and re-validation policy.]
    The spec does not yet define how \texttt{ValidatedPlan} values are
    persisted, versioned, and invalidated in the web platform's plan store
    (Blob Storage, Cosmos DB).  Phase~2 must specify retention, eviction,
    and forced re-validation triggers (e.g., grammar upgrade, runbook source
    change, admin override).

  \item[\texttt{OPQ-GATE-06} — Conformance corpus coverage for gate-level
    codes.]
    Stream~C (Tess) must add test vectors for \texttt{PLAN-001} through
    \texttt{PLAN-009} to the conformance corpus.  Stream~F does not define
    these vectors; they are a Phase~2 corpus deliverable.  The gate-level
    codes that require migration-era negative vectors are \texttt{PLAN-004}
    (Go-template detection), \texttt{PLAN-005} (forbidden infix), and
    \texttt{PLAN-006} (pipe capture).

  \item[\texttt{OPQ-GATE-07} — OI-GXL-01 through OI-GXL-07 and
    OI-GIS-02 through OI-GIS-05 (grammar open issues).]
    The 12~open issues from \texttt{gxl.ebnf~\S7} and
    \texttt{gis.ebnf~\S6} that were not resolved in Stream~A are queued for
    Phase~2.  The parse gate is not directly affected by them, but the
    conformance corpus and evaluator implementations must account for them.
    The most gate-relevant are:
    \begin{itemize}
      \item \texttt{OI-GXL-04} (\texttt{not null} behaviour / \texttt{defined()}
        builtin request): may produce a new \texttt{GXL-PARSE-011} in v2.
      \item \texttt{OI-GIS-04} (nested \texttt{\$\{...\}} inside GXL string
        literals): requires parser coordination not yet specified.
      \item \texttt{OI-GCP-03} (statically reachable step set, also
        \texttt{OPQ-GATE-02} above).
      \item \texttt{OI-GCP-05} (\texttt{http.*} scope lifetime): affects
        \texttt{PLAN-009} reachability semantics.
    \end{itemize}
\end{description}
